\rtmtight
\begin{tblr}{width=\textwidth,colspec={|Q[c,m,2.1cm]|X[l,m]|},hlines,vlines,hspan=minimal,rowsep=1.5pt,colsep=3pt,row{1}={bg=headgray,font=\bfseries}}
\SetCell{c,m}\hfil\logo\hfil & \textbf{POLITEKNIK NEGERI MALANG}\par\textbf{JURUSAN TEKNOLOGI INFORMASI}\par\textbf{PROGRAM STUDI: D4 TEKNIK INFORMATIKA} \\
\SetCell[c=2]{l} \textbf{RENCANA TUGAS MAHASISWA (RTM) DAN RUBRIK PENILAIAN} & \\
Nama Mata Kuliah & Administrasi dan Keamanan Jaringan \\
Kode Mata Kuliah & RTI255008 \quad \textbf{sks / Jam perminggu:} 3 SKS/6 jam \quad \textbf{Semester:} 5 \\
Dosen Pengampu & \begin{enumerate}\item Ade Ismail\item Yuri Ariyanto\item Sofyan Noor Arief\item Yusriel A\end{enumerate} \\
\end{tblr}
\assessmentblock{Tugas Analisis Desain Keamanan Jaringan}{Tugas terstruktur dan studi kasus}{SCPMK0501-04401, SCPMK0502-04401, SCPMK0505-04401, SCPMK0505-04402}{Mahasiswa menyelesaikan tugas analisis desain keamanan jaringan sesuai Sub-CPMK dan konteks mata kuliah, mencakup identifikasi ancaman, perancangan topologi aman, konfigurasi administrasi, dan optimasi jaringan.}{Menganalisis skenario keamanan jaringan, merancang solusi teknis, mengkonfigurasi mekanisme keamanan, mendokumentasikan evidence, dan mempresentasikan keputusan bila diminta.}{Laporan teknis, diagram topologi jaringan, konfigurasi terdokumentasi, evidence pengujian, dan/atau bahan presentasi.}{Ketepatan identifikasi ancaman dan kerentanan jaringan & 9 & Ancaman, kerentanan, dan dampak diidentifikasi secara komprehensif dengan contoh nyata dan referensi SKKNI atau best practice keamanan jaringan. & Ancaman utama teridentifikasi tetapi analisis dampak atau referensi standar belum lengkap. & Identifikasi ancaman tidak tepat atau tidak mencerminkan konteks keamanan jaringan yang nyata. \\ \\
Kualitas desain dan konfigurasi keamanan jaringan & 9 & Desain mencakup segmentasi, DMZ, penempatan perangkat keamanan, dan konfigurasi yang diuji serta didokumentasikan secara lengkap. & Desain cukup valid tetapi konfigurasi, pengujian, atau justifikasi komponen belum sepenuhnya lengkap. & Desain tidak menunjukkan prinsip keamanan jaringan yang benar atau konfigurasi tidak berfungsi. \\ \\
Kelengkapan dokumentasi dan laporan teknis & 6 & Dokumentasi lengkap, sistematis, mencakup desain, konfigurasi, pengujian, dan refleksi teknis yang siap digunakan sebagai referensi. & Dokumentasi ada tetapi belum sistematis atau masih ada bagian penting yang kurang. & Dokumentasi tidak lengkap atau tidak dapat digunakan sebagai referensi teknis. \\}

\assessmentblock{Simulasi Desain Keamanan Jaringan}{Project Based Learning -- Praktikum/Simulasi}{SCPMK0501-04401, SCPMK0502-04401}{Mahasiswa menyelesaikan simulasi desain keamanan jaringan sesuai Sub-CPMK, mencakup implementasi topologi aman, konfigurasi firewall, dan dokumentasi hasil simulasi menggunakan GNS3 atau Packet Tracer.}{Menganalisis skenario, merancang dan mengimplementasikan topologi jaringan aman, mengkonfigurasi aturan firewall dan penempatan perangkat keamanan, menguji konektivitas dan keamanan, serta mendokumentasikan laporan proyek.}{Laporan proyek, diagram topologi terdokumentasi, konfigurasi dan evidence pengujian, dan/atau bahan presentasi.}{Akurasi implementasi desain keamanan jaringan & 5 & Topologi, segmentasi, DMZ, dan penempatan perangkat keamanan diimplementasikan sesuai skenario dan best practice dengan bukti pengujian yang lengkap. & Implementasi utama benar tetapi beberapa detail segmentasi, penempatan perangkat, atau pengujian belum optimal. & Implementasi tidak mencerminkan prinsip desain keamanan jaringan atau tidak dapat diuji. \\ \\
Konfigurasi firewall dan aturan akses & 5 & Aturan firewall dikonfigurasi dengan logika keamanan yang benar, diuji dengan skenario serangan, dan seluruh hasilnya terdokumentasi. & Konfigurasi dasar ada dan sebagian pengujian dilakukan, tetapi dokumentasi aturan atau skenario pengujian belum lengkap. & Konfigurasi tidak berfungsi atau tidak mengikuti logika keamanan yang benar. \\ \\
Kualitas laporan dan presentasi simulasi & 3 & Laporan mencakup desain, implementasi, hasil pengujian, dan refleksi dengan argumen teknis yang kuat dan sistematis. & Laporan ada tetapi analisis teknis, kelengkapan bagian, atau keterkaitan antar komponen masih kurang. & Laporan tidak sistematis atau tidak mendokumentasikan proses simulasi secara memadai. \\}

\assessmentblock{Kuis Konfigurasi dan Keamanan Jaringan}{Kuis tertulis}{SCPMK0501-04401, SCPMK0505-04401}{Mahasiswa mengerjakan kuis tertulis yang mencakup pemahaman konsep konfigurasi gateway, NAT, firewall, dan mekanisme keamanan jaringan sesuai Sub-CPMK.}{Menjawab soal kuis secara tertulis berdasarkan pemahaman konsep dan kemampuan analisis skenario konfigurasi jaringan.}{Lembar jawaban kuis tertulis.}{Penguasaan konsep konfigurasi gateway dan NAT & 4 & Jawaban menunjukkan pemahaman mendalam tentang konfigurasi NAT, firewall, routing, dan keamanan gateway berdasarkan skenario yang diberikan. & Konsep dasar dipahami tetapi penerapan pada skenario spesifik masih kurang tepat atau tidak lengkap. & Jawaban tidak menunjukkan pemahaman konsep konfigurasi jaringan yang memadai. \\ \\
Analisis skenario keamanan jaringan & 4 & Skenario dianalisis dengan tepat, solusi yang diusulkan sesuai best practice, dan dampak keamanan diidentifikasi dengan benar. & Analisis cukup baik tetapi rekomendasi solusi atau identifikasi dampak belum sepenuhnya lengkap. & Analisis tidak relevan atau solusi yang diusulkan tidak mengikuti prinsip keamanan jaringan. \\ \\
Ketepatan prosedur implementasi keamanan & 2 & Prosedur implementasi dikuasai dengan urutan dan justifikasi yang benar sesuai standar industri. & Prosedur sebagian benar tetapi ada langkah yang salah, terlewat, atau tidak dapat dijustifikasi. & Prosedur tidak dikuasai, urutannya salah, atau tidak sesuai standar implementasi keamanan jaringan. \\}

\assessmentblock{UTS Infrastruktur dan Keamanan Jaringan}{Ujian Tertulis}{SCPMK0501-04401, SCPMK0502-04401}{Mahasiswa mengerjakan ujian tertulis komprehensif yang mencakup materi infrastruktur jaringan, identifikasi ancaman, penilaian risiko, desain topologi aman, konsep gateway, dan konfigurasi NAT/firewall dari pertemuan 1--7.}{Menjawab soal ujian tertulis secara mandiri yang mencakup konsep, analisis skenario, dan perancangan solusi keamanan jaringan.}{Lembar jawaban ujian tertulis.}{Pemahaman konsep infrastruktur dan keamanan jaringan & 8 & Mahasiswa mampu menjelaskan, menganalisis, dan menghubungkan konsep CIA Triad, ancaman, kerentanan, kebijakan keamanan, dan desain infrastruktur secara komprehensif. & Konsep utama dipahami tetapi hubungan antar konsep atau penerapan pada skenario masih kurang lengkap. & Jawaban tidak menunjukkan pemahaman konsep infrastruktur dan keamanan jaringan yang memadai. \\ \\
Analisis skenario ancaman dan kerentanan & 7 & Ancaman dan kerentanan diidentifikasi dengan tepat, dampak bisnis dianalisis, dan strategi mitigasi yang sesuai diusulkan berdasarkan best practice. & Identifikasi ancaman cukup baik tetapi analisis dampak bisnis atau strategi mitigasi belum lengkap. & Ancaman tidak teridentifikasi dengan tepat atau mitigasi yang diusulkan tidak relevan dengan konteks. \\ \\
Desain solusi keamanan jaringan & 5 & Solusi dirancang dengan mempertimbangkan topologi, segmentasi, penempatan perangkat, dan mekanisme perlindungan yang tepat dan terjustifikasi. & Solusi cukup valid tetapi belum mencakup semua aspek keamanan yang relevan atau justifikasi masih lemah. & Solusi tidak menunjukkan pemahaman desain keamanan jaringan atau tidak dapat diterapkan. \\}

\assessmentblock{Simulasi Administrasi dan Troubleshooting Jaringan}{Project Based Learning -- Praktikum/Simulasi}{SCPMK0505-04401, SCPMK0505-04402}{Mahasiswa menyelesaikan simulasi administrasi dan troubleshooting jaringan sesuai Sub-CPMK, mencakup konfigurasi administrasi akses, manajemen pengguna, monitoring SIEM, dan penyelesaian skenario troubleshooting.}{Menganalisis skenario administrasi jaringan, mengkonfigurasi akses aman (SSH, AAA, RBAC), mensimulasikan skenario troubleshooting, menggunakan alat monitoring, dan mendokumentasikan seluruh proses.}{Laporan proyek, konfigurasi terdokumentasi, log dan evidence monitoring, bahan troubleshooting, dan/atau bahan presentasi.}{Konfigurasi administrasi dan kontrol akses jaringan & 5 & Administrasi akses (SSH, MFA, AAA), manajemen pengguna (RBAC, Least Privilege), dan konfigurasi log berhasil dikonfigurasi, diuji, dan terdokumentasi secara lengkap. & Konfigurasi utama berhasil tetapi beberapa fitur administrasi, pengujian, atau dokumentasi belum lengkap. & Konfigurasi tidak berfungsi atau tidak sesuai dengan prinsip keamanan administrasi jaringan. \\ \\
Analisis dan troubleshooting masalah jaringan & 5 & Masalah jaringan diidentifikasi dengan tepat, didiagnosis secara sistematis menggunakan alat yang sesuai, dan diselesaikan dengan solusi yang terdokumentasi. & Identifikasi masalah cukup baik tetapi langkah troubleshooting atau dokumentasi solusi belum lengkap. & Troubleshooting tidak sistematis, masalah tidak dapat diselesaikan, atau tidak ada dokumentasi. \\ \\
Penggunaan monitoring dan logging (SIEM) & 3 & Alat monitoring dikonfigurasi, log dianalisis untuk keperluan keamanan, dan insight dari SIEM digunakan untuk meningkatkan keamanan sistem. & Monitoring dan logging dilakukan tetapi analisis log atau integrasi SIEM masih terbatas. & Tidak ada pemahaman atau penerapan yang memadai tentang monitoring dan logging jaringan. \\}

\assessmentblock{UAS Kompetensi Administrasi dan Keamanan Jaringan}{Ujian Akhir -- Tertulis atau Presentasi Proyek}{SCPMK0501-04401, SCPMK0502-04401, SCPMK0505-04401, SCPMK0505-04402}{Mahasiswa mengerjakan ujian akhir atau mempresentasikan proyek akhir yang menunjukkan kompetensi komprehensif dalam administrasi dan keamanan jaringan sesuai SKKNI dan capaian seluruh CPMK mata kuliah.}{Mengerjakan soal ujian tulis atau mempresentasikan proyek akhir yang mencakup desain, implementasi, administrasi, monitoring, optimasi, dan evaluasi sistem keamanan jaringan secara menyeluruh.}{Lembar jawaban ujian atau laporan proyek akhir lengkap dengan dokumentasi teknis, hasil pengujian, dan bahan presentasi.}{Kompetensi administrasi jaringan terpadu & 8 & Mahasiswa menunjukkan kemampuan mengelola, mengkonfigurasi, dan memelihara sistem jaringan secara komprehensif sesuai SKKNI dengan bukti teknis yang meyakinkan. & Kompetensi administrasi cukup baik tetapi beberapa aspek pengelolaan atau konfigurasi lanjutan belum sepenuhnya dikuasai. & Tidak dapat menunjukkan kompetensi administrasi jaringan yang memadai atau bukti teknis tidak meyakinkan. \\ \\
Strategi keamanan dan perlindungan sistem informasi & 7 & Strategi keamanan komprehensif dirancang dan diimplementasikan, mencakup identifikasi, autentikasi, otorisasi, monitoring, dan kebijakan perlindungan CIA Triad. & Strategi keamanan mencakup aspek utama tetapi beberapa mekanisme perlindungan atau integrasi antar komponen belum optimal. & Strategi keamanan tidak lengkap atau tidak menunjukkan pemahaman konsep perlindungan sistem informasi. \\ \\
Optimasi kinerja dan presentasi kompetensi akhir & 5 & Hasil optimasi kinerja jaringan dipresentasikan dengan evidence data, analisis bottleneck, rekomendasi yang dapat ditindaklanjuti, dan komunikasi teknis yang meyakinkan. & Presentasi cukup baik tetapi analisis data kinerja, rekomendasi, atau kemampuan mempertahankan argumen masih terbatas. & Tidak dapat mempresentasikan hasil kompetensi dengan bukti teknis yang memadai atau komunikasi tidak profesional. \\}

\begin{tblr}{width=\textwidth,colspec={|Q[l,m,3.2cm]|X[l,m]|},hlines,vlines,hspan=minimal,row{1-3}={bg=headgray},rowsep=1.5pt,colsep=3pt}
Jadwal Pelaksanaan: & Minggu 1--17 sesuai RPS. Setiap evaluasi memiliki RTM dan rubrik multi-kriteria. \\
Lain-Lain Yang Diperlukan: & LMS, simulator jaringan (GNS3/VirtualBox), repositori konfigurasi, dan perangkat presentasi. \\
Pustaka: & Mengacu pustaka utama dan pendukung pada bagian RPS. \\
\end{tblr}
